\subsection{Microsoft Access-এ ডেটাবেজ ও টেবিল তৈরির কাজের ধাপ}
ডেটাবেজ তৈরিতে প্রস্তুত Template/Wizard থেকে শুরু করা, Blank database নিয়ে নিজে নকশা করা অথবা সমর্থিত ভাষার নির্দেশ দিয়ে কাঠামো তৈরি করা যায়। Template-এর তৈরি ফিল্ড ও নিয়মও নিজের প্রয়োজনের সঙ্গে মিলিয়ে নিতে হয়। সব সফটওয়্যারে একই Wizard বা \texttt{CREATE DATABASE} নির্দেশ থাকে না; SQLite-এ নতুন ফাইল খুলে তার মধ্যে টেবিল তৈরি করা হয়।

Access খুলে Blank database নির্বাচন করে ফাইলের নাম ও সংরক্ষণের স্থান নির্ধারণ করি এবং Create দিই। আধুনিক Access ফাইলের প্রচলিত এক্সটেনশন \texttt{.accdb}; পুরোনো ফরম্যাটে \texttt{.mdb} ব্যবহৃত হয়। ফাইলের নাম ও তার ভেতরের টেবিলের নাম এক জিনিস নয়। একটি College.accdb ফাইলে Student, Result ও Payment টেবিল থাকতে পারে।

পরের Member উদাহরণে একই কাজের ধাপ অনুসরণ করা হয়েছে: Blank Database → Table Design → Field Name → Data Type → Primary Key → Field Properties → Save → Datasheet View → নমুনা রেকর্ড যাচাই। Student টেবিল তৈরি করলেও একই ধাপ প্রযোজ্য; সেখানে StudentID-এর সারিতে Primary Key দিতে হবে। AutoNumber সাধারণত সফটওয়্যার তৈরি করে; অন্য ফিল্ডের মান ব্যবহারকারী লেখেন।

ফিল্ড ও টাইপের পরিকল্পনা আগে করলে অপ্রয়োজনীয় পরিবর্তন কমে। যেমন Mobile-কে Number করে শুরুতে শূন্য হারানোর পর শুধু প্রদর্শনের বিন্যাস বদলে মূল পরিচিতি নির্ভরযোগ্যভাবে ফিরে পাওয়া যায় না। তাই প্রথমেই Text নির্বাচন করা ভালো। বিস্তারিত টাইপ ও বৈশিষ্ট্য পরের অংশে আছে।

\chSixInterfaceSide{ch06_access_create_database.png}{Access-এ নতুন ডেটাবেজ তৈরির পর্দার পরিচিতি}{চিত্রের Blank database নির্বাচন করে File Name ঘরে ডেটাবেজের নাম দিতে হয়। Create বোতাম নতুন ফাইল তৈরি করে; এর ভেতরে পরে একাধিক টেবিল রাখা যায়।}
\chSixInterfaceSide{ch06_access_table_design.png}{Table Design-এ ফিল্ড, ডেটা টাইপ ও প্রাইমারি কী}{উপরের প্রতিটি সারি একটি ফিল্ডের কাঠামো বোঝায়। নিচের Field Properties নির্বাচিত ফিল্ডের নিয়ম নির্ধারণ করে; কী চিহ্নটি প্রাইমারি কী নির্দেশ করে।}

\subsubsection{Member টেবিলে কাজের ধাপ-এর প্রয়োগ}
\textbf{১. ফিল্ডের নাম:} প্রথম সারিতে MemberID, পরের সারিতে Name, JoinDate ও Mobile লিখি। প্রতিটি নাম একটি বৈশিষ্ট্য বোঝাবে। একই টেবিলে একই নামের দুটি ফিল্ড দেওয়া যায় না। একটি দীর্ঘ ঠিকানার সব অংশ একই ফিল্ডে থাকবে, নাকি District ও Area আলাদা হবে, তা প্রয়োজনীয় অনুসন্ধান দেখে নির্ধারণ করি।

\textbf{২. টাইপ ও বিবরণ:} MemberID-এ AutoNumber, Name ও Mobile-এ Short Text এবং JoinDate-এ Date/Time নিই। Description-এ ফিল্ডের উদ্দেশ্য লেখা যায়; যেমন Mobile-এ ``যোগাযোগের নম্বর''। এই বর্ণনা ফিল্ডের ডেটা টাইপ বা সংরক্ষিত মানের বিকল্প নয়।

\textbf{৩. কী ও প্রপার্টি:} MemberID-এর সারি নির্বাচন করে Primary Key দিই। Name প্রয়োজনীয় হলে Required=Yes এবং প্রযোজ্য Text সেটিংয়ে Allow Zero Length=No নির্ধারণ করি। জন্ম বা যোগদানের তারিখের মতো ফিল্ডে গ্রহণযোগ্য সীমা প্রতিষ্ঠানের নিয়ম অনুযায়ী দিই। Type নির্বাচনের পরেই তার উপযোগী property দেখা যায়।

\textbf{৪. সংরক্ষণ:} টেবিল Member নামে Save করি। এরপর Datasheet View-তে যাই। ফাইলের নাম Library.accdb হলেও ভেতরের টেবিলের নাম Member। একটি ডেটাবেজে পরে Issue ও Book টেবিলও তৈরি করা যাবে। তাই টেবিলের নাম ও ডেটাবেজ ফাইলের নাম আলাদা উদ্দেশ্য প্রকাশ করে।

\textbf{৫. রেকর্ড এন্ট্রি ও যাচাই:} তিন সদস্যের তথ্য লিখে মোট তিনটি ডেটা-সারি আছে কি না দেখি। শিরোনামের সারি গুনব না। MemberID স্বয়ংক্রিয়ভাবে তৈরি হয়েছে কি না এবং Mobile-এর শুরুতে শূন্য আছে কি না যাচাই করি। একজন সদস্যের নাম সংশোধন করে পুনরায় খোলার পর সঠিক মান আছে কি না দেখি।

